\newpage
\section{Net sensitivity: a constraint group that cannot break what it did not look at}

\subsection*{Shift factor, revisited}

Every quantity a constraint group is built from is already in hand: row $e$ of
the sensitivity matrix (\S\ref{app:dump}, Stage 4) gives, for branch $e$, the
change in $|f_e|$ per MW curtailed at every wind node --- signed so that
positive means \emph{relief}. [WDT]'s method (Appendix 2, p.\,72) reads a
\emph{single} row: pick the one congested circuit, rank nodes down that row,
keep whatever clears a threshold $\tau$.
\begin{equation}
  \text{group}(e) \;=\; \{\, n : \text{sensitivity}_{e,n} \,\ge\, \tau \,\},
  \qquad \tau = \pTau .
\end{equation}
That is one row of a $\gne\times\gnwind$ matrix, chosen and read in isolation.
Nothing in it depends on any other row --- so nothing in it \emph{knows} about
any other row.

\subsection*{Sensitivities add. Net sensitivity is nothing more than that.}

If several branches are over their rating at once, a node's total usefulness is
just its sensitivities to each of them, summed:
\begin{equation}
  S_{\text{net}}(n) \;=\; \sum_{e\,\in\,\text{Over}} \text{sensitivity}_{e,n},
  \qquad \text{Over} = \{\, e : |f_e| > s_e \,\}.
  \label{eq:netsens}
\end{equation}
No new computation --- the matrix already holds every term on the right. A
node relieving three circuits a little now outranks one relieving only the
worst circuit a lot; a one-line group is unaffected, since \eqref{eq:netsens}
collapses to the single term $\tau$ was already thresholding.

\subsection*{Why net sensitivity alone is not enough}

Ranking by $S_{\text{net}}$ still asks only \emph{"does this node help the
branches I am watching."} It says nothing about the $\gne - |\text{Over}|$
branches it is not watching. Tested directly: ranking by $S_{\text{net}}$ with
no further check pushed a branch that was inside its rating over it in
\pOvlBrokeEvents\ of \pN\ events --- \emph{worse} than the threshold method it
replaces. Concentrating a whole instruction on the single most useful node
moves the rest of the network harder than the same megawatts spread out.
\textbf{The safety was never going to come from the ranking.} It has to come
from a separate rule.

\subsection*{The veto}

A node's cut is capped, not by $\tau$, but by the largest amount that leaves
\emph{every live branch} $k$ --- not just the ones in Over --- inside its
rating, and leaves nothing already over made worse:
\begin{equation}
  x_n^{\max} \;=\; \min_{k \,:\, \text{sensitivity}_{k,n} < 0}
  \;\frac{\text{rating}_k - |f_k|}{|\text{sensitivity}_{k,n}|},
  \qquad
  \text{lim}_k = \max\!\big(|f_k|,\ \text{rating}_k\big) \text{ for } k \in \text{Over}.
  \label{eq:veto}
\end{equation}
Every branch in the network enters this minimum, whether it is congested or
not --- because the branch a bad instruction breaks is, by definition, one
that was not being watched. A node is included in the group only up to
$x_n^{\max}$; ordering by $S_{\text{net}}$ decides \emph{who goes first},
\eqref{eq:veto} decides \emph{how far they may go}. The two are independent
checks and the results below keep them separate.

\begin{center}\small
\begin{tabular}{lrrrr}
\toprule
& \textbf{cleared} & \textbf{events broke a line} & \textbf{lines broken} & \textbf{median group}\\
\midrule
Shift-factor threshold (today) & \pEirCleared\% & \pEirBrokeEvents & \pEirBrokeLines & \pEirMedNodes\ nodes\\
Net sensitivity + veto & \pVetoCleared\% & \pVetoBrokeEvents & \pVetoBrokeLines & \pVetoMedNodes\ node\\
\bottomrule
\end{tabular}
\end{center}
\pN\ (outage, hour) events, every single-branch outage at every hour of the
week. On the events both methods clear, the veto method cuts
\pSavedMW\ MW less wind --- \pSavedPct\% --- because it targets the effective
node instead of splitting pro-rata across \pEirMedNodes. The proven optimum
(the same rule in \eqref{eq:veto}, solved exactly rather than applied greedily)
clears \pLpCleared\%; a safe remedy provably exists in \pLpFeasible\% of events.

\subsection*{Fixed, and precomputable, per outage}

Neither \eqref{eq:netsens} nor \eqref{eq:veto} depends on the hour or the
weather --- both are read off the sensitivity matrix, which is topology and
reactance only. So the group for \emph{every} single-branch outage can be
built once, offline, and looked up: \pTripsOutages\ outages produce
\pTripsCircuits\ distinct circuits ever at risk, \pTrips\ (outage, circuit)
lookup rows in total. \pVetoedFromEir\ of the \pTripsCircuits-circuit
threshold method's own \pEirMedNodes-node-median memberships
(\pVetoedFromEirPct\% of every member ever proposed) fail \eqref{eq:veto} and
are excluded before dispatch. Only the live flow and the final MW instruction
are computed online.

\subsection*{Finding 1 --- curtailment is concentrated on a handful of nodes}

Running \eqref{eq:veto} across every one of the \pN\ events and tracking which
nodes the safe remedy actually dispatches, not merely which ones qualify as
candidates: no node is cut in \emph{zero} events --- the least-used one is
still asked \pMinCutEvents\ times out of \pN --- so there is no node that can
be declared exempt outright. What the same sweep shows instead is
concentration. The \textbf{5} busiest nodes account for \pTopFivePct\% of every
megawatt curtailed anywhere in the study; the top 10 for \pTopTenPct\%; the top
20 for \pTopTwentyPct\%. \pRareN\ of \pWindN\ nodes --- \pRarePct\% of the fleet,
\pRareMW\ MW of nameplate, \pRareFleetPct\% of the \pWindFleetMW\ MW installed
--- are cut in under 1\% of the \pN\ events, and the median node across the
whole fleet is cut in only \pMedEventsPct\% of them, keeping \pMedKeptPct\% of
whatever it was offering when it was. Overall the safe remedy cuts
\pOverallCutPct\% of all wind offered across the entire N--1 study:
\pTotalCutMW\ MW of \pTotalOfferedMW\ MW. \textbf{For the great majority of
the fleet, a constraint group is something that applies to them on a handful
of days a year, not a standing limit on their output.}

\subsection*{Finding 2 --- stopping the cascade}

Follow the chain past the first trip: whatever is left over its rating after
the remedy trips too, and the flows are recomputed with it also out of
service. \pCascN\ starting outages, run to exhaustion under four remedies:

\begin{center}\small
\begin{tabular}{lrrr}
\toprule
& \textbf{lines lost, total} & \textbf{outages that spread} & \textbf{islanded}\\
\midrule
No remedy & \pNothingLost & \pNothingSpread\ / \pCascN & \pNothingIsl\\
Shift-factor threshold & \pCEirLost & \pCEirSpread\ / \pCascN & \pCEirIsl\\
Net sensitivity + veto & \pCVetoLost & \pCVetoSpread\ / \pCascN & \pCVetoIsl\\
Safe optimum (LP) & \pCLpLost & \pCLpSpread\ / \pCascN & \pCLpIsl\\
\bottomrule
\end{tabular}
\end{center}
\note{This sample is deliberately enriched with the events where a remedy was
already seen to break something, so the spread rate above is not the rate for
a uniformly random N--1 --- it is a stress test, by design.}

\medskip
\note{\textbf{Limits.} No thermal clock: a branch over its rating at the end
of a round is taken to trip, the pessimistic reading, since real protection is
time-overcurrent. The veto is exact only inside the DC linearisation it shares
with the rest of this document. Weather is synthetic, so the event count and
the MW figures move with real data; \emph{which} circuits are ever at risk does
not, since it follows from $K$, $B$ and the ratings alone.}
